%
% mpv-program-player.tex
%
% Z Specification for the voxd program tier's mpv PROCESS/CONNECTION lifecycle.
%
% This model is the stateful-audio gate for docs/mpv-program-player.md. It
% models the lifecycle of the ONE persistent mpv process and its JSON-IPC
% connection --- spawn, connect, crash-detect, restart, give-up, shutdown ---
% and the loop's interaction with it (wait-for-ready, loadfile, receive
% end-file, replay-after-crash).
%
% It is DISTINCT from, and does not restate, the domain source state machine of
% docs/vox-music-player-transport.tex. That model (T1--T7: idle/playing/paused,
% the cursor, the catalogue) is UNCHANGED by the mpv mechanism swap; the design
% doc's Section 3 argues T1--T7 hold because they are properties of the source
% state, which mpv does not touch. This model pins the new stateful behaviour
% the swap introduces: the process/connection lifecycle and its eight
% invariants I1--I7 and single-loadfile-ownership.
%
% It reuses the style, naming, and the ZBOOL convention of the transport model,
% and the same bounded-integer discipline so ProB can exhaust the reachable
% space.
%

\documentclass[a4paper,10pt,fleqn]{article}
\usepackage[margin=1in]{geometry}
\usepackage{fuzz}
\usepackage[colorlinks=true,linkcolor=blue,citecolor=blue,urlcolor=blue]{hyperref}

\begin{document}

\title{The voxd Program Player: A Z Specification of the mpv Process and
Connection Lifecycle}
\author{Formal model of the mpv process/connection state machine}
\date{August 2026}
\hypersetup{
  pdftitle={The voxd Program Player: A Z Specification of the mpv Process and Connection Lifecycle},
  pdfauthor={Punt Labs},
  pdfsubject={Z Specification},
  pdfcreator={fuzz/probcli}
}
\maketitle

\section{Introduction}

The program tier abandons a fresh \texttt{ffplay} process per part in favour of
\emph{one} persistent \texttt{mpv} process driven over a JSON IPC socket
(\texttt{docs/mpv-program-player.md}). Playing a part becomes \texttt{loadfile};
a part ending becomes an \texttt{end-file} event; pause becomes an internal
property flip, click-free because nothing is torn down. The change swaps the
\emph{mechanism} of playback and leaves the \emph{domain} state machine ---
the source cursor and the transport modes \texttt{idle}/\texttt{playing}/%
\texttt{paused} of \texttt{docs/vox-music-player-transport.tex}, with its
invariants \texttt{T1}--\texttt{T7} --- untouched. That model is not restated
here, exactly as the transport model did not restate the phase-2 catalogue
operations it depended on.

What the swap \emph{introduces} is a new, genuinely stateful subsystem: the
lifecycle of the one mpv process and its socket connection, and the discipline
by which the playback loop drives it without ever issuing a command into a dead
or not-yet-ready connection, without losing its place across a crash, and
without racing a second loader against the supervisor. This specification pins
that lifecycle. It fixes the process modes, the terms the design names
(the pending-command count, the loop's in-flight await, the loadfile-issuer,
the restart counter, the paused flag, and the current-part cursor), and the
eight invariants \texttt{I1}--\texttt{I7} and \emph{single-loadfile-ownership}
that must hold across every transition.

The division of labour is the settled one (\texttt{md}~\S1, peer-reviewed
\texttt{gvr}+\texttt{kpz}): a \emph{supervisor} owns the process and the
connection --- spawn, connect, crash-detect, restart --- and \emph{nothing}
about what plays; the \emph{loop} owns every \texttt{loadfile} and holds the
cursor. This model makes that split a theorem: \texttt{loadfileBy} never
contains the supervisor, and no \texttt{loadfile} is enabled outside
\texttt{ready}, reached only by \emph{waiting} for \texttt{ready}
(\texttt{WaitReady}), never by issuing-and-retrying.

\section{Basic Types, Modes, and Constants}
\label{sec:types}

The mpv process is in one of six modes over its lifecycle.

\begin{zed}
  MpvState ::= down | starting | ready | crashed | restarting | failed
\end{zed}

\begin{description}
  \item[\texttt{down}] no process --- before the first spawn, and after a
    clean shutdown.
  \item[\texttt{starting}] the process is spawned but the socket is not yet
    connected --- the first-time bring-up.
  \item[\texttt{ready}] the process is alive, the socket is connected and its
    reader is running, and commands are accepted. This is the \emph{only} mode
    in which any IPC command may be issued.
  \item[\texttt{crashed}] the process or socket has died; on entry the reader
    resolves every pending command future and the loop's in-flight ended-future
    (with the synthetic reason \texttt{crashed}), and a restart is owed.
  \item[\texttt{restarting}] the process is being respawned after a crash,
    within the cap; like \texttt{starting} but carrying the standing fault the
    crash raised.
  \item[\texttt{failed}] the restart cap is exceeded; the supervisor has given
    up, and a standing hard fault stands until the daemon is shut down.
\end{description}

The playback loop is in one of three phases with respect to the mpv connection.
The phase is what makes \texttt{WaitReady} an explicit, reachable step rather
than an implicit precondition.

\begin{zed}
  LoopPhase ::= lwait | lready | lplay
\end{zed}

\begin{description}
  \item[\texttt{lwait}] the loop is parked in \texttt{WaitReady}, awaiting the
    \texttt{ready} signal before it will issue any \texttt{loadfile} --- the
    startup wait, and the post-crash recovery wait.
  \item[\texttt{lready}] the loop has observed \texttt{ready} (via
    \texttt{WaitReady}) and is cleared to issue its \texttt{loadfile}; reached
    \emph{only} from \texttt{lwait} when the process is \texttt{ready}, so the
    path to a \texttt{loadfile} is provably a wait, not a retry.
  \item[\texttt{lplay}] a file is loaded and the loop is awaiting its
    ended-future (\texttt{loopAwait} = \texttt{ztrue}).
\end{description}

An \texttt{end-file} carries one of five reasons. The first four are mpv's own;
\texttt{rCrashed} is \emph{synthetic} --- the reader injects it into the loop's
ended-future on socket EOF, because a crash emits no \texttt{end-file}. It is
produced by \texttt{Crash} (Section~\ref{sec:crash}), never by a real event.

\begin{zed}
  EndFileReason ::= rEof | rStop | rRedirect | rError | rCrashed
\end{zed}

The two actors that could conceivably issue a \texttt{loadfile}: the loop and
the supervisor. The single-loadfile-ownership invariant restricts the issuers
to $\{ loopA \}$.

\begin{zed}
  ACTOR ::= loopA | supA
\end{zed}

A \texttt{ZBOOL} is a Z boolean, named to avoid the ProB/B clash with
\texttt{BOOL}, exactly as in the transport model. It flags the standing fault,
the loop's in-flight await, the paused flag, and the load's pause option.

\begin{zed}
  ZBOOL ::= ztrue | zfalse
\end{zed}

Four bounds keep the carriers finite so ProB can explore the reachable space.
\texttt{maxRestarts} is the restart cap; \texttt{maxCmds} bounds the in-flight
command futures; \texttt{maxPart} bounds the part cursor as in the transport
model. None is a product limit; each is a model bound chosen small enough to
exhaust and large enough to exercise every boundary --- the cap reached, more
than one command in flight, and a cursor that can rotate.

\begin{axdef}
  maxRestarts : \nat_1 \\
  maxCmds : \nat_1 \\
  maxPart : \nat_1
  \where
  maxRestarts = 2 \\
  maxCmds = 2 \\
  maxPart = 2
\end{axdef}

\section{State}
\label{sec:state}

The \texttt{MpvLifecycle} carries the process mode; the number of live
processes and readers; the restart counter; a standing-fault flag; the count of
unresolved command futures; the loop phase; the loop's in-flight-await flag;
the set of actors that have issued the current \texttt{loadfile}; the paused
flag; the pause option the current file was loaded with; and the current-part
cursor.

Four of these are \emph{derived}: \texttt{processes}, \texttt{readers}, and
\texttt{fault} are total functions of \texttt{state}, and \texttt{loopAwait} is
a total function of \texttt{loopPhase}. They are declared and pinned in the
predicate --- as the transport model declared and pinned \texttt{advancing} and
\texttt{glyph} --- so the invariants they carry (\texttt{I2}, \texttt{I5},
\texttt{I3}, and the await-side of \texttt{I7}) read as facts in the state, not
as prose. The predicate carries every invariant; read the conjuncts in the
order \texttt{I1}--\texttt{I7} then ownership.

\begin{schema}{MpvLifecycle}
  state : MpvState \\
  processes : 0 \upto 1 \\
  readers : 0 \upto 1 \\
  restarts : 0 \upto maxRestarts \\
  fault : ZBOOL \\
  pendingCmds : 0 \upto maxCmds \\
  loopPhase : LoopPhase \\
  loopAwait : ZBOOL \\
  loadfileBy : \power ACTOR \\
  paused : ZBOOL \\
  loadPause : ZBOOL \\
  cursor : 1 \upto maxPart
  \where
  (loopPhase \in \{ lready, lplay \} \implies state = ready) \\
  (state \neq ready \implies pendingCmds = 0) \\
  (processes = 1 \iff state \in \{ starting, restarting, ready \}) \\
  (readers = 1 \iff state = ready) \\
  (state = ready \iff processes = 1 \land readers = 1) \\
  (fault = ztrue \iff state \in \{ crashed, restarting, failed \}) \\
  (state = failed \implies restarts = maxRestarts) \\
  (loopAwait = ztrue \iff loopPhase = lplay) \\
  loadfileBy \subseteq \{ loopA \}
\end{schema}

The conjuncts discharge the invariants as follows.

\begin{description}
  \item[\texttt{I1} --- commands only when \texttt{ready}, reached by waiting.]
    The state face is $loopPhase \in \{ lready, lplay \} \implies state =
    ready$ together with $state \neq ready \implies pendingCmds = 0$: no
    command future stands, and the loop has neither cleared-to-issue nor an
    active load, in any non-\texttt{ready} mode. The \emph{operational} face
    --- that a \texttt{loadfile} is reached by \emph{waiting} on \texttt{ready}
    and not by a refuse-and-retry --- is carried by \texttt{WaitReady}
    (Section~\ref{sec:waitready}): \texttt{lready} is set \emph{only} there, and
    \emph{only} from \texttt{lwait} when $state = ready$. So \texttt{I1} is
    proven by the wait, not discharged vacuously.
  \item[\texttt{I2} --- at most one process.] $processes \in 0 \upto 1$ by
    type, pinned to $1$ exactly in \texttt{starting}/\texttt{restarting}/%
    \texttt{ready}. No restart race can spawn a second, because \texttt{Spawn}
    fires only from \texttt{down} and \texttt{Restart} only from
    \texttt{crashed}, and neither of those has a live process.
  \item[\texttt{I3} --- fault iff \texttt{crashed}/\texttt{restarting}/%
    \texttt{failed}.] $fault = ztrue \iff state \in \{ crashed, restarting,
    failed \}$. This is the schema-sketch reading of \texttt{I3}, which is
    \emph{stricter} than the prose reading ``fault iff not ready'' and,
    Section~\ref{sec:forks} argues, the correct one: \texttt{down} and
    \texttt{starting} are non-\texttt{ready} but carry no standing fault, being
    clean shutdown/pre-spawn and clean first bring-up.
  \item[\texttt{I4} --- the restart cap terminates.] $restarts \in 0 \upto
    maxRestarts$ by type, and $state = failed \implies restarts =
    maxRestarts$. \texttt{Restart} increments \texttt{restarts} and is enabled
    only while $restarts < maxRestarts$; at the cap only \texttt{Fail} is
    enabled, and \texttt{failed} spawns nothing. There is no hot respawn loop.
  \item[\texttt{I5} --- one reader.] $readers \in 0 \upto 1$ by type, pinned
    to $1$ exactly in \texttt{ready}: the reader is alive iff the connection is
    up, started on \texttt{Connect} and gone on \texttt{Crash}/\texttt{Shutdown}.
  \item[\texttt{I6} --- the lifecycle does not corrupt the cursor, and recovery
    honours pause.] A property of the crash/recovery operations, not a standing
    conjunct: \texttt{Crash}, \texttt{Restart}, and \texttt{ReplayCurrent} all
    hold $cursor' = cursor$ (Sections~\ref{sec:crash}, \ref{sec:replay}), and
    \texttt{ReplayCurrent} sets $loadPause' = paused$ --- the reload is loaded
    paused exactly when the paused flag stands. Recovery reaches its
    \texttt{loadfile} through \texttt{WaitReady}, never a busy-retry.
  \item[\texttt{I7} --- no orphaned await across a crash.] The await face is
    $loopAwait = ztrue \iff loopPhase = lplay$, and $loopPhase \in \{ lready,
    lplay \} \implies state = ready$; the command face is $state \neq ready
    \implies pendingCmds = 0$. Together: entering any non-\texttt{ready} mode
    forces $pendingCmds = 0$ and $loopAwait = zfalse$. \texttt{Crash}
    (Section~\ref{sec:crash}) realises this: it sets $pendingCmds' = 0$ (fails
    every pending command future) and $loopPhase' = lwait$ (resolves the
    ended-future, with the synthetic \texttt{crashed} reason). No await is left
    hanging.
  \item[\emph{single-loadfile-ownership}.] $loadfileBy \subseteq \{ loopA \}$:
    only the loop ever issues \texttt{loadfile}. \texttt{Spawn},
    \texttt{Connect}, \texttt{Crash}, \texttt{Restart}, \texttt{Fail},
    \texttt{CrashDuringRestart}, and \texttt{Shutdown} --- the supervisor's
    transitions --- never write \texttt{loadfileBy}; only
    \texttt{ReplayCurrent} and \texttt{AutoAdvance} do, and both set it to
    $\{ loopA \}$. The supervisor never enters the set.
\end{description}

Because \texttt{processes}, \texttt{readers}, and \texttt{fault} are total
functions of \texttt{state}, and \texttt{loopAwait} of \texttt{loopPhase},
every state fixes them uniquely; they add no freedom, only a name for a fact.
The genuinely free variables are \texttt{state}, \texttt{restarts},
\texttt{pendingCmds}, \texttt{loopPhase}, \texttt{loadfileBy}, \texttt{paused},
\texttt{loadPause}, and \texttt{cursor}, and the operations below frame all
eight; the four derived observations follow from the state invariant of
\texttt{MpvLifecycle$'$} and are not written out, exactly as
\texttt{advancing$'$}/\texttt{glyph$'$} were left to the invariant in the
transport model.

\section{Initialisation}

The daemon starts \texttt{down}: no process, the loop parked in
\texttt{WaitReady}, the restart counter and command futures at zero, nothing
paused. The boot cursor is left free over $1 \upto maxPart$ so the animator
explores every starting part; an initialisation establishes a state, it does
not consume an input.

\begin{schema}{MpvInit}
  MpvLifecycle'
  \where
  state' = down \\
  restarts' = 0 \\
  pendingCmds' = 0 \\
  loopPhase' = lwait \\
  loadfileBy' = \emptyset \\
  paused' = zfalse \\
  loadPause' = zfalse
\end{schema}

From $state' = down$ and $loopPhase' = lwait$ the invariant forces
$processes' = 0$, $readers' = 0$, $fault' = zfalse$, and $loopAwait' =
zfalse$; only \texttt{cursor$'$} is free.

\section{Transitions}

Each operation names its enabling mode on an early predicate line, then frames
the eight free variables. The supervisor owns \texttt{Spawn}, \texttt{Connect},
\texttt{Crash}, \texttt{Restart}, \texttt{Fail}, \texttt{CrashDuringRestart},
and \texttt{Shutdown}; the loop owns \texttt{WaitReady}, \texttt{ReplayCurrent},
\texttt{AutoAdvance}, and \texttt{EofWhilePaused}; the service owns
\texttt{SetPause}, \texttt{SetResume}, and \texttt{ResolveCommand}.

\subsection{Spawn --- bring the process up}
\label{sec:spawn}

From \texttt{down}, \texttt{Spawn} starts the process; the socket is not yet
connected, so the mode is \texttt{starting}, not \texttt{ready}.

\begin{schema}{Spawn}
  \Delta MpvLifecycle
  \where
  state = down \\
  state' = starting \\
  restarts' = restarts \\
  pendingCmds' = pendingCmds \\
  loopPhase' = loopPhase \\
  loadfileBy' = loadfileBy \\
  paused' = paused \\
  loadPause' = loadPause \\
  cursor' = cursor
\end{schema}

The invariant forces $processes' = 1$, $readers' = 0$, $fault' = zfalse$: one
process, no reader, no fault --- a clean bring-up.

\subsection{Connect --- the connection comes up}
\label{sec:connect}

From \texttt{starting} (first bring-up) or \texttt{restarting} (post-crash
respawn), \texttt{Connect} attaches the reader and reaches \texttt{ready}. The
two are one path, as the design intends: ``first connect'' and ``reconnect
after crash'' run the identical routine. A successful reconnect resets the
restart counter --- this is how the model renders the design's ``restart cap
within a window'': the counter measures \emph{consecutive} respawns that never
reached \texttt{ready}, and reaching \texttt{ready} clears it.

\begin{schema}{Connect}
  \Delta MpvLifecycle
  \where
  state \in \{ starting, restarting \} \\
  state' = ready \\
  restarts' = 0 \\
  pendingCmds' = pendingCmds \\
  loopPhase' = loopPhase \\
  loadfileBy' = loadfileBy \\
  paused' = paused \\
  loadPause' = loadPause \\
  cursor' = cursor
\end{schema}

The invariant forces $readers' = 1$ and $fault' = zfalse$. Because
$pendingCmds = 0$ held in \texttt{starting}/\texttt{restarting}, $pendingCmds'
= 0$ still satisfies the ready-side of the state predicate, which permits any
count. The loop phase is unchanged --- typically \texttt{lwait}, waiting for
exactly this signal.

\subsection{WaitReady --- the loop waits, it does not retry}
\label{sec:waitready}

From \texttt{lwait}, and only once the process is \texttt{ready}, the loop
transitions to \texttt{lready}: cleared to issue its \texttt{loadfile}. This is
the whole content of \texttt{I1}'s operational face. The loop parks in
\texttt{lwait} exactly as long as mpv takes to come up, and no longer; it never
issues a \texttt{loadfile} into a not-ready client and gets refused into a
fault, then spins.

\begin{schema}{WaitReady}
  \Delta MpvLifecycle
  \where
  state = ready \\
  loopPhase = lwait \\
  state' = state \\
  restarts' = restarts \\
  pendingCmds' = pendingCmds \\
  loopPhase' = lready \\
  loadfileBy' = loadfileBy \\
  paused' = paused \\
  loadPause' = loadPause \\
  cursor' = cursor
\end{schema}

\subsection{ReplayCurrent --- (re)play the current part, honouring pause}
\label{sec:replay}

From \texttt{lready}, the loop issues the \texttt{loadfile} for the
\emph{current} part. This is both the first play at startup and the replay
after a crash --- one operation, because the design unifies them. The cursor is
unmoved (\texttt{I6}), the loader is the loop (\emph{ownership}), and the load's
pause option equals the paused flag (\texttt{I6}: a crash while paused replays
paused, so recovery never silently resumes playing). The loop enters
\texttt{lplay}, now awaiting the ended-future.

\begin{schema}{ReplayCurrent}
  \Delta MpvLifecycle
  \where
  state = ready \\
  loopPhase = lready \\
  state' = state \\
  restarts' = restarts \\
  pendingCmds' = pendingCmds \\
  loopPhase' = lplay \\
  loadfileBy' = \{ loopA \} \\
  paused' = paused \\
  loadPause' = paused \\
  cursor' = cursor
\end{schema}

The invariant forces $loopAwait' = ztrue$ from $loopPhase' = lplay$. Because
$loadfileBy' = \{ loopA \}$, ownership is preserved; because $loadPause' =
paused$, the reload honours the paused flag whatever its value; because
$cursor' = cursor$, the part does not move.

\subsection{AutoAdvance --- a part ends cleanly, advance and reload}
\label{sec:autoadvance}

While \texttt{lplay} and \emph{not paused}, an \texttt{end-file} with reason
\texttt{rEof} (natural end) or \texttt{rError} (bad part, advance past it)
rotates the cursor to another part of the album --- avoiding an immediate
repeat when the album holds more than one part, replaying when it holds exactly
one --- and immediately reloads, staying in \texttt{lplay}. The paused guard is
the retained \texttt{T3} check (\texttt{loop.py:159}): an \texttt{rEof} that mpv
buffered in the instant before the user paused must not advance-then-load while
the mode reads paused.

\begin{schema}{AutoAdvance}
  \Delta MpvLifecycle \\
  reason? : EndFileReason
  \where
  state = ready \\
  loopPhase = lplay \\
  reason? \in \{ rEof, rError \} \\
  paused = zfalse \\
  state' = state \\
  restarts' = restarts \\
  pendingCmds' = pendingCmds \\
  loopPhase' = lplay \\
  loadfileBy' = \{ loopA \} \\
  paused' = paused \\
  loadPause' = paused \\
  1 \leq cursor' \land cursor' \leq maxPart \\
  (maxPart \geq 2 \implies cursor' \neq cursor) \\
  (maxPart = 1 \implies cursor' = cursor)
\end{schema}

Only the cursor moves, and it stays in $1 \upto maxPart$. The load is the
loop's ($loadfileBy' = \{ loopA \}$) and is played, not paused ($loadPause' =
paused = zfalse$), so ownership and \texttt{I6} hold.

\subsection{EofWhilePaused --- the guarded \texttt{rEof} is dropped}
\label{sec:eofpaused}

While \texttt{lplay} and \emph{paused}, an \texttt{end-file} with reason
\texttt{rEof} is dropped: the cursor does not move and the loop keeps awaiting.
This is the other face of the \texttt{T3} guard. A paused mpv reaches no
\emph{new} \texttt{rEof}, but the one it buffered just before the pause must not
advance the cursor; the loop ignores it and stays where it is.

\begin{schema}{EofWhilePaused}
  \Xi MpvLifecycle \\
  reason? : EndFileReason
  \where
  state = ready \\
  loopPhase = lplay \\
  reason? = rEof \\
  paused = ztrue
\end{schema}

$\Xi MpvLifecycle$ holds every variable fixed: the guarded event changes
nothing.

\subsection{Crash --- the connection dies}
\label{sec:crash}

From \texttt{ready}, socket EOF (or a set \texttt{returncode}) is a crash. On
detection the reader fails every pending command future and resolves the loop's
in-flight ended-future with the synthetic reason \texttt{crashed}; the model
renders both as $pendingCmds' = 0$ and $loopPhase' = lwait$ (which forces
$loopAwait' = zfalse$). The cursor is untouched and the paused flag survives ---
a crash while paused stays paused. The current load is gone, so $loadfileBy' =
\emptyset$.

\begin{schema}{Crash}
  \Delta MpvLifecycle
  \where
  state = ready \\
  state' = crashed \\
  restarts' = restarts \\
  pendingCmds' = 0 \\
  loopPhase' = lwait \\
  loadfileBy' = \emptyset \\
  paused' = paused \\
  loadPause' = loadPause \\
  cursor' = cursor
\end{schema}

The invariant forces $processes' = 0$, $readers' = 0$, $fault' = ztrue$. This
is the \texttt{I7} transition: no await --- neither a command future nor the
ended-future --- survives the crash. \texttt{Crash} is enabled from every
\texttt{ready} state, whatever the loop phase; when the loop was not in
\texttt{lplay} there is no ended-future to resolve, but the pending command
futures are failed and the loop is pushed to \texttt{lwait} to recover all the
same.

\subsection{Restart --- respawn within the cap}
\label{sec:restart}

From \texttt{crashed}, while the cap is not yet reached, \texttt{Restart}
increments the counter and respawns, reaching \texttt{restarting}. The standing
fault persists (\texttt{restarting} carries it) until \texttt{Connect} clears
it.

\begin{schema}{Restart}
  \Delta MpvLifecycle
  \where
  state = crashed \\
  restarts < maxRestarts \\
  state' = restarting \\
  restarts' = restarts + 1 \\
  pendingCmds' = pendingCmds \\
  loopPhase' = loopPhase \\
  loadfileBy' = loadfileBy \\
  paused' = paused \\
  loadPause' = loadPause \\
  cursor' = cursor
\end{schema}

The invariant forces $processes' = 1$ (respawned), $readers' = 0$ (not yet
connected), $fault' = ztrue$.

\subsection{CrashDuringRestart --- the respawn dies before connecting}
\label{sec:crashrestart}

A respawn that dies before it reaches \texttt{ready} returns to \texttt{crashed}
without resetting the counter. This is the path that lets the counter climb to
the cap: a run of respawns that never connect drives \texttt{restarts} to
\texttt{maxRestarts}, at which point only \texttt{Fail} is enabled.

\begin{schema}{CrashDuringRestart}
  \Delta MpvLifecycle
  \where
  state = restarting \\
  state' = crashed \\
  restarts' = restarts \\
  pendingCmds' = pendingCmds \\
  loopPhase' = loopPhase \\
  loadfileBy' = loadfileBy \\
  paused' = paused \\
  loadPause' = loadPause \\
  cursor' = cursor
\end{schema}

\subsection{Fail --- the cap is reached, give up}
\label{sec:fail}

From \texttt{crashed} at the cap, \texttt{Fail} enters \texttt{failed}: a
standing hard fault, no further respawn. This is \texttt{I4}'s terminating
edge.

\begin{schema}{Fail}
  \Delta MpvLifecycle
  \where
  state = crashed \\
  restarts = maxRestarts \\
  state' = failed \\
  restarts' = restarts \\
  pendingCmds' = pendingCmds \\
  loopPhase' = loopPhase \\
  loadfileBy' = loadfileBy \\
  paused' = paused \\
  loadPause' = loadPause \\
  cursor' = cursor
\end{schema}

The after-state satisfies $state = failed \implies restarts = maxRestarts$; the
invariant forces $processes' = 0$, $fault' = ztrue$. No transition from
\texttt{failed} respawns; only \texttt{Shutdown} leaves it.

\subsection{SetPause, SetResume --- direct pause/resume}
\label{sec:pause}

The service pauses and resumes by a direct IPC command in \texttt{ready}
(settled: not routed through the control channel, because they mutate no source
state). Each registers a command future --- captured here as an increment of
\texttt{pendingCmds} --- and flips the paused flag. Neither reloads: the load's
pause option \texttt{loadPause} is untouched, which is the click-free property.

\begin{schema}{SetPause}
  \Delta MpvLifecycle
  \where
  state = ready \\
  pendingCmds < maxCmds \\
  state' = state \\
  restarts' = restarts \\
  pendingCmds' = pendingCmds + 1 \\
  loopPhase' = loopPhase \\
  loadfileBy' = loadfileBy \\
  paused' = ztrue \\
  loadPause' = loadPause \\
  cursor' = cursor
\end{schema}

\begin{schema}{SetResume}
  \Delta MpvLifecycle
  \where
  state = ready \\
  pendingCmds < maxCmds \\
  state' = state \\
  restarts' = restarts \\
  pendingCmds' = pendingCmds + 1 \\
  loopPhase' = loopPhase \\
  loadfileBy' = loadfileBy \\
  paused' = zfalse \\
  loadPause' = loadPause \\
  cursor' = cursor
\end{schema}

Both require \texttt{ready} --- part of \texttt{I1}: no command is issued into
a not-ready connection.

\subsection{ResolveCommand --- a response arrives}
\label{sec:resolve}

A command's response arrives and its future resolves, decrementing the
in-flight count. This is the ordinary counterpart to \texttt{Crash}, which
resolves them all at once with an error.

\begin{schema}{ResolveCommand}
  \Delta MpvLifecycle
  \where
  state = ready \\
  pendingCmds > 0 \\
  state' = state \\
  restarts' = restarts \\
  pendingCmds' = pendingCmds - 1 \\
  loopPhase' = loopPhase \\
  loadfileBy' = loadfileBy \\
  paused' = paused \\
  loadPause' = loadPause \\
  cursor' = cursor
\end{schema}

\subsection{Shutdown --- graceful teardown}
\label{sec:shutdown}

From any live mode, \texttt{Shutdown} quits mpv and returns to \texttt{down}:
the process and reader gone, the counter and futures cleared, the loop parked
to wait again should the daemon come back up. Enabled from \texttt{failed} as
well, so a failed process is not a wedge --- the daemon can be shut down and a
fresh \texttt{Spawn} started.

\begin{schema}{Shutdown}
  \Delta MpvLifecycle
  \where
  state \neq down \\
  state' = down \\
  restarts' = 0 \\
  pendingCmds' = 0 \\
  loopPhase' = lwait \\
  loadfileBy' = \emptyset \\
  paused' = zfalse \\
  loadPause' = zfalse \\
  cursor' = cursor
\end{schema}

The invariant forces $processes' = 0$, $readers' = 0$, $fault' = zfalse$: a
clean \texttt{down}, indistinguishable from the pre-spawn state save for the
cursor it leaves in place.

\section{Design forks and observations surfaced by the formalisation}
\label{sec:forks}

Formalising the lifecycle exposes one genuine fork and three modelling
decisions worth recording for the leader. The model takes a stated position in
each and flags the tension rather than papering over it.

\begin{description}
  \item[Fork~1 --- the reading of \texttt{I3}. \emph{Recommend the strict
    reading.}] The design states \texttt{I3} two ways that do not agree. The
    prose (\S6, invariant list) says ``fault $\iff$ not ready''; the state-schema
    sketch (\S6) says $fault = ztrue \iff state \in \{ crashed, restarting,
    failed \}$. These differ on \texttt{down} and \texttt{starting}, which are
    non-\texttt{ready} but, under the sketch, fault-free. The two clean
    non-\texttt{ready} modes --- \texttt{down} (pre-spawn or post-clean-shutdown)
    and \texttt{starting} (normal first bring-up) --- carry no standing fault:
    marking them faulted would report a fault during every ordinary startup and
    after every ordinary shutdown, which is not what a client should see.
    \textbf{This model takes the strict schema reading}, $fault \iff state \in
    \{ crashed, restarting, failed \}$, and recommends the design prose be
    corrected to match. \emph{Decision needed from the leader before
    implementation:} confirm the strict reading, so \texttt{doctor} and
    \texttt{ProgramStatus.playback\_error} report a standing fault in
    \texttt{crashed}/\texttt{restarting}/\texttt{failed} only, and treat the
    spawn-failure \texttt{PLAYER\_UNAVAILABLE} case as the next item rather than
    as ``\texttt{starting} with a fault''.
  \item[Observation~A --- \texttt{PLAYER\_UNAVAILABLE} is outside this model.]
    The design's spawn-failure fault (mpv missing or too old at bring-up) is a
    \emph{startup} edge, not part of the crash/recovery lifecycle the brief's
    transitions cover; there is no \texttt{SpawnFailure} transition in the set.
    It surfaces as a standing \texttt{PLAYER\_UNAVAILABLE} fault with the daemon
    still up, and it does \emph{not} fit the six-mode/\texttt{I3} shape cleanly
    (a failed spawn has $restarts = 0$, so it is not \texttt{failed}, and it
    carries a fault, so it is not \texttt{down}/\texttt{starting} under the
    strict \texttt{I3}). \textbf{Recommendation:} model it, if at all, as a
    seventh mode \texttt{unavailable} distinct from \texttt{failed}, or accept
    that the startup-availability fault lives outside this lifecycle's
    invariant. It is bracketed here as scope, exactly as the transport model
    bracketed the catalogue operations. Flagged because it is the one place the
    strict \texttt{I3} and the design's fault surface must be reconciled in
    code.
  \item[Observation~B --- \texttt{restarting} versus reusing \texttt{starting}.]
    The design sketch's \texttt{Restart} reads ``\texttt{crashed} $\to$
    \texttt{starting}''. This model routes the post-crash respawn through the
    distinct \texttt{restarting} mode instead, so that \texttt{I3}'s fault
    biconditional stays clean: \texttt{restarting} carries the standing fault
    the crash raised, whereas \texttt{starting} (first bring-up) does not. The
    two modes are otherwise identical --- process up, socket not yet connected
    --- and \texttt{Connect} serves both. This is a refinement, not a
    contradiction; recorded so the implementation keeps the standing fault up
    across a respawn until \texttt{ready} is re-reached.
  \item[Observation~C --- the restart window is modelled as consecutive failed
    respawns.] The design specifies a ``restart cap within a window''. Modelling
    wall-clock windows is out of a Z model's reach; this model renders the
    window as \emph{consecutive respawns that never reached \texttt{ready}}:
    \texttt{restarts} climbs on each \texttt{Restart} and \texttt{Connect}
    resets it to $0$. A single successful reconnect clears the count. This is a
    faithful and checkable abstraction of the intent --- ``do not hot-loop'' ---
    and the implementation is free to add a time bound on top; the invariant it
    must preserve, \texttt{I4}, is unaffected.
\end{description}

\section{Type-checking and model-checking}
\label{sec:results}

The specification is \texttt{fuzz}-clean: \texttt{fuzz -t} exits~0.

It was then model-checked with ProB (\texttt{probcli}, $maxRestarts = 2$,
$maxCmds = 2$, $maxPart = 2$, integers bounded to $0 \upto 3$). Over the boot
cursor ranging across every part, ProB visits every reachable state
(125~states, 485~transitions) and finds \emph{no invariant violation and no
deadlock}: the \texttt{MpvLifecycle} predicate --- and with it \texttt{I1}
(the state face), \texttt{I2}, \texttt{I3}, \texttt{I4}, \texttt{I5}, the await
side of \texttt{I7}, and single-loadfile-ownership --- holds on every reachable
state, and every state has an enabled successor. All fourteen operations are
covered, so each is reachable and exercised. The absence of a deadlock is the
key liveness result: a \texttt{crashed} state is never a wedge, because
\texttt{Restart}/\texttt{Fail} always leave it and \texttt{Shutdown} leaves any
live mode.

The invariants that are properties of the crash/recovery \emph{operations}
rather than standing conjuncts were confirmed by temporal checks over the full
reachable space:

\begin{description}
  \item[\texttt{I7} --- no orphaned await.] $\mathrm{AG}(state = crashed
    \implies loopAwait = zfalse \land pendingCmds = 0)$ holds: every crash
    resolves both the loop's ended-future and every pending command future.
    That the loop \emph{genuinely} awaits and that \emph{several} futures can be
    in flight --- so this is not vacuous --- is witnessed by
    $\mathrm{EF}(loopAwait = ztrue)$ and $\mathrm{EF}(pendingCmds = maxCmds)$,
    both reachable.
  \item[Crash liveness.] $\mathrm{AG}(state = crashed \implies \mathrm{EF}(state
    = ready \lor state = failed))$ holds: from any crashed state the machine can
    always reach \texttt{ready} again or terminate at \texttt{failed} --- it
    never wedges.
  \item[\texttt{I4} --- the cap terminates.] $\mathrm{EF}(state = failed)$ is
    reachable, and $\mathrm{EF}(state = failed \land restarts = maxRestarts)$
    confirms \texttt{failed} is entered only at the cap; the state invariant
    $state = failed \implies restarts = maxRestarts$ holds throughout. No
    reachable state respawns from \texttt{failed}.
  \item[\texttt{I6} --- recovery honours pause.] $\mathrm{EF}(loopPhase = lplay
    \land loadPause = ztrue \land paused = ztrue)$ is reachable: a witness of
    length~12 drives the machine to pause a playing part, crash, restart,
    reconnect, wait, and replay --- and the replayed load carries
    $loadPause = ztrue$ because the paused flag stood. A crash while paused
    does not silently resume playing.
  \item[single-loadfile-ownership.] $\mathrm{AG}(supA \notin loadfileBy)$ holds:
    across the whole reachable space the supervisor never appears among the
    \texttt{loadfile} issuers. The double-load the earlier supervisor-reloads
    split allowed is foreclosed.
  \item[\texttt{I1} --- commands only when ready.] $\mathrm{AG}(state \neq ready
    \implies pendingCmds = 0)$ holds, and \texttt{lready} --- the loop's
    cleared-to-issue phase --- is set only by \texttt{WaitReady} from
    \texttt{lwait} at $state = ready$, so the path to every \texttt{loadfile} is
    a wait on \texttt{ready}, not a refuse-and-retry.
\end{description}

The model is deadlock-free by construction of \texttt{Shutdown} (enabled from
any live mode) and \texttt{Spawn} (from \texttt{down}); unlike the transport
model, there is no scope-boundary dead end, because the process lifecycle is
self-contained --- it does not depend on an unmodelled catalogue operation to
make progress.

\end{document}
